\section{Background and Related Work}
\label{sec:related-work}

\subsection{Reproducibility and experimental variability in machine learning}
\label{sec:related-reproducibility}

Concerns about reproducibility in artificial intelligence and machine learning
span both research reporting and experimental practice. Gundersen and Kjensmo
surveyed empirical AI publications and found substantial gaps in the
documentation needed to reproduce reported experiments
\cite{gundersen2018reproducibility}. Raff subsequently studied independent
reimplementation of machine-learning research and emphasized that the
availability of code alone does not determine whether a result can be
independently reproduced \cite{raff2019reproducible}. Community interventions
such as the NeurIPS reproducibility program have therefore combined artifact
availability with reporting checklists and other changes to experimental
practice \cite{pineau2021improving}.

A related body of work has shown that apparently minor experimental choices can
materially affect measured ML performance. Reimers and Gurevych demonstrated
that random-seed choice can change reported results for neural sequence-tagging
systems \cite{reimers2017score}. Henderson et al.\ documented substantial
variance and reporting challenges in deep reinforcement learning
\cite{henderson2018matters}. Bouthillier et al.\ analyzed variance arising from
data sampling, parameter initialization, and hyperparameter selection in ML
benchmarks \cite{bouthillier2021variance}, while Pham et al.\ showed that both
algorithmic and implementation-level factors can introduce variance in deep
learning software \cite{pham2020variance}.

Reproducibility problems in ML-based science are not limited to stochastic
variation. Kapoor and Narayanan documented data leakage as a recurring source
of overoptimistic scientific conclusions \cite{kapoor2023leakage}. Heil et
al.\ similarly argue for explicit reproducibility standards spanning data,
models, code, programming practice, and workflow automation
\cite{heil2021standards}.

These studies motivate treating experimental configuration as part of the
scientific artifact rather than as incidental software detail. However, their
primary questions concern reproducibility, reporting, variance, or
methodological correctness. Our study asks a complementary testing question:
when a reproducibility-relevant experimental choice is deliberately changed,
does a validation workflow already present in the repository detect that
change?


\subsection{Executability and quality of research software}
\label{sec:related-executability}

Artifact availability is distinct from artifact executability. Collberg and
Proebsting reported substantial difficulties obtaining and rebuilding software
artifacts from computer-systems research \cite{collberg2016repeatability}.
At a larger scale, Trisovic et al.\ studied thousands of research-code files
from publicly available replication datasets and found that a large fraction
failed during clean-environment re-execution, even though some failures could
be mitigated through code cleaning \cite{trisovic2022execution}.

This distinction matters for mutation-based empirical studies. A mutation
cannot be meaningfully classified as detected or undetected when the
unmodified baseline cannot first execute under the selected validation
workflow. For this reason, our protocol treats setup, environment, and baseline
workflow failures as a separate evidence layer rather than assigning them
mutation outcomes.

The same literature also motivates preserving environment and workflow
information as part of research artifacts. Reproducibility recommendations for
ML emphasize explicit software environments, workflow automation, and
publication of the materials necessary to rerun an analysis
\cite{heil2021standards}. In our study, these concerns arise not only as
recommendations for future projects but also as practical constraints on
retrospective execution of historical repositories.


\subsection{Mutation testing and scientific software}
\label{sec:related-mutation}

Mutation testing evaluates a validation mechanism by introducing controlled
changes---mutants---and observing whether the available tests distinguish the
mutated program from the original. The technique has a long history in
software engineering; Jia and Harman provide a broad survey of its development,
operator design, cost, and mutation adequacy measures
\cite{jia2011mutation}.

A central justification for mutation analysis is that artificial faults can
provide a useful proxy for assessing test effectiveness when the complete set
of real faults is unknown. Just et al.\ compared mutant detection with
detection of real faults and found empirical evidence supporting this
relationship while also identifying limitations of the proxy
\cite{just2014mutants}. Mutation studies must also contend with mutants that do
not change relevant program semantics, motivating explicit treatment of
equivalent mutants rather than counting every syntactic change as an effective
fault.

Mutation techniques have also been considered in scientific software, where
the oracle problem can be particularly acute. Hook and Kelly used mutation in
the context of testing trustworthiness of scientific software, explicitly
addressing the difficulty of determining correct numerical outputs
\cite{hook2009trustworthiness}. This work is important precedent for applying
mutation concepts outside conventional business software.

Our use of mutation testing follows the same general adequacy principle but
changes the object being measured. We do not primarily ask whether a unit-test
suite detects arbitrary program faults. Instead, we use narrowly defined
reproducibility-relevant transformations to probe whether a frozen validation
workflow constrains experimental choices that may influence a research result.


\subsection{Mutation testing of machine-learning systems}
\label{sec:related-ml-mutation}

Mutation testing is already well established as a research direction for
machine-learning and deep-learning systems. DeepMutation introduced both
source-level mutations of training data and training programs and model-level
mutations of trained neural networks, using mutant detection to assess test-data
quality \cite{ma2018deepmutation}. DeepMutation++ extended mutation-based
evaluation for deep neural networks, including support for feed-forward and
recurrent architectures \cite{hu2019deepmutationpp}.

Subsequent work moved toward mutation operators grounded in observed DL faults.
DeepCrime derives source-level mutation operators from real fault evidence and
uses them to assess deep-learning test data
\cite{humbatova2021deepcrime}. DeepMetis in turn uses mutation score as an
optimization objective when augmenting DL test sets
\cite{riccio2021deepmetis}. These studies demonstrate that domain-specific
mutation operators can be more informative than simply applying conventional
program mutations to ML code.

There is therefore no novelty claim in this paper that mutation testing has not
previously been applied to ML systems. Some prior source-level DL mutation
operators also touch training and experimental configuration. The distinction
in our work is the measurement target and the oracle: MLReproMutate targets
reproducibility-relevant choices in research repositories and evaluates them
against validation workflows that existed in those repositories, rather than
primarily assessing the adequacy of model test data or the robustness of a
trained neural network.


\subsection{Dependencies, environments, and workflow reproducibility}
\label{sec:related-environments}

Dependency changes form an especially important overlap between reproducibility
engineering and mutation testing. A declared dependency modification is not
necessarily equivalent to a semantic environment change: package resolution
may select the same installed version, or a changed version may affect only
code that is not exercised by validation.

Hejderup and Gousios directly studied whether test suites can be trusted to
support automated dependency updates. Using Java projects and a large number of
artificial dependency faults, they found substantial gaps in the ability of
tests to detect semantic problems introduced through dependencies
\cite{hejderup2022dependencies}. Their work demonstrates that mutation-based
evaluation of dependency-update safeguards predates the present study.

Our \texttt{dependency-pin} operator addresses a different but adjacent
question in research software. It changes a repository's declared exact
constraint from \texttt{package==version} to
\texttt{package>=version}, resolves baseline and mutant environments
independently, and verifies whether the installed target dependency actually
changed before interpreting the validation result. The selected repository
workflow, rather than dependency-call coverage alone, serves as the empirical
oracle.

Environment concerns also connect dependency management to long-term
executability. Research-code studies show that software and environment drift
can prevent later execution \cite{trisovic2022execution}, while reproducibility
standards emphasize capturing and automating computational workflows
\cite{heil2021standards}. This motivates our separation between semantic
dependency changes, validation outcomes, and infrastructure failures.


\subsection{Positioning of this study}
\label{sec:related-positioning}

The literature above establishes substantial prior work on ML reproducibility,
research-code executability, conventional and scientific-software mutation
testing, mutation testing of deep-learning systems, and testing of dependency
updates. The contribution of the present study is therefore not a claim to be
the first use of mutation testing in ML, scientific software, or dependency
management.

Instead, the study combines these concerns around a specific empirical
measurement target: the sensitivity of existing ML research-repository
validation workflows to controlled changes in reproducibility-relevant
experimental choices. The empirical protocol operates on frozen revisions of
real research repositories, freezes mutation candidates and validation
workflows before observing outcomes, distinguishes infrastructure failures from
mutation outcomes, and explicitly verifies semantic equivalence where required.

This positioning also distinguishes mutation survival from a direct measure of
reproducibility. A surviving confirmed non-equivalent mutation shows that the
selected frozen workflow did not detect one controlled change. It does not
establish that the repository or its reported scientific result is
irreproducible. Conversely, a killed mutation demonstrates sensitivity to that
specific change rather than general reproducibility assurance.
